\section{Conclusion}
This study introduced Adaptive VST-PGD, a transparent low-dose PET preprocessing pipeline that combines Poisson variance stabilization, spatially adaptive first- and second-order regularization, and projected spectral-gradient optimization. Across three simulated dose levels, the method achieved the highest PSNR and competitive edge preservation, while NLM produced higher SSIM. Ablation results identified variance stabilization as the dominant component and confirmed a measurable contribution from adaptive weighting.

The numerical model omits scatter, random events, attenuation, resolution loss, reconstruction-induced correlation, and heterogeneous clinical anatomy. Periodic boundary operators may also introduce border interactions. Future work should evaluate reflective operators, the exact unbiased Anscombe inverse, direct Poisson-likelihood optimization, accelerated solvers, automated parameter selection, and lesion-level task metrics. Validation should progress from multiple digital phantoms to physical phantoms and ethically approved patient cohorts before clinical conclusions are drawn.
